%!TEX TS-program = lualatex
%!TEX encoding = UTF-8 Unicode
%!BIB program = bibtex

\documentclass{mimosis}

\usepackage{metalogo}

%%%%%%%%%%%%%%%%%%%%%%%%%%%%%%%%%%%%%%%%%%%%%%%%%%%%%%%%%%%%%%%%%%%%%%%%
% Some of my favourite personal adjustments
%%%%%%%%%%%%%%%%%%%%%%%%%%%%%%%%%%%%%%%%%%%%%%%%%%%%%%%%%%%%%%%%%%%%%%%%
%
% These are the adjustments that I consider necessary for typesetting
% a nice thesis. However, they are *not* included in the template, as
% I do not want to force you to use them.

% This ensures that I am able to typeset bold font in table while still aligning the numbers
% correctly.
\usepackage{etoolbox}

\usepackage{siunitx}
\DeclareSIUnit\px{px}

\sisetup{%
  detect-all           = true,
  detect-family        = true,
  detect-mode          = true,
  detect-shape         = true,
  detect-weight        = true,
}

%%%%%%%%%%%%%%%%%%%%%%%%%%%%%%%%%%%%%%%%%%%%%%%%%%%%%%%%%%%%%%%%%%%%%%%%
% Hyperlinks & bookmarks
%%%%%%%%%%%%%%%%%%%%%%%%%%%%%%%%%%%%%%%%%%%%%%%%%%%%%%%%%%%%%%%%%%%%%%%%

\usepackage[%
  unicode,
  colorlinks = true,
  citecolor  = RoyalBlue,
  linkcolor  = RoyalBlue,
  urlcolor   = RoyalBlue,
  unicode,
  ]{hyperref}
\hypersetup{
  pdftitle=On Language Modeling and its Applications for Contemporary and Historical French,
  pdfauthor=Pedro Ortiz Suarez
}

\usepackage{bookmark}


%%%%%%%%%%%%%%%%%%%%%%%%%%%%%%%%%%%%%%%%%%%%%%%%%%%%%%%%%%%%%%%%%%%%%%%%
% Fonts
%%%%%%%%%%%%%%%%%%%%%%%%%%%%%%%%%%%%%%%%%%%%%%%%%%%%%%%%%%%%%%%%%%%%%%%%

\ifxetexorluatex
  \setmainfont[Ligatures=TeX]{TeX Gyre Pagella} % Palatino clone
  \linespread{1.05} % a bit more for Palatino
  \RequirePackage{unicode-math}
  \setmathfont{TeX Gyre Pagella Math}
\else
  \usepackage[lf]{ebgaramond}
  \usepackage[oldstyle,scale=0.7]{sourcecodepro}
  \singlespacing
\fi

\usepackage{microtype}

\renewcommand{\th}{\textsuperscript{\textup{th}}\xspace}

\newacronym[description={Principal component analysis}]{PCA}{PCA}{principal component analysis}
\newacronym                                            {SNF}{SNF}{Smith normal form}
\newacronym[description={Topological data analysis}]   {TDA}{TDA}{topological data analysis}

\newglossaryentry{LaTeX}{%
  name        = {\LaTeX},
  description = {A document preparation system},
  sort        = {LaTeX},
}

\newglossaryentry{Real numbers}{%
  name        = {$\real$},
  description = {The set of real numbers},
  sort        = {Real numbers},
}

\makeindex
\makeglossaries

%%%%%%%%%%%%%%%%%%%%%%%%%%%%%%%%%%%%%%%%%%%%%%%%%%%%%%%%%%%%%%%%%%%%%%%%
% Custom Commands and Environments
%%%%%%%%%%%%%%%%%%%%%%%%%%%%%%%%%%%%%%%%%%%%%%%%%%%%%%%%%%%%%%%%%%%%%%%%

%\usepackage{fontawesome}
\usepackage{arydshln}
\usepackage{fontawesome5}
\usepackage{tikz}
\usetikzlibrary{positioning}

\usepackage{stackengine}
\def\ucr{\scalebox{1}{\stackinset{c}{}{c}{-.1pt}{%
  \textcolor{white}{\sffamily\bfseries\small ?}}{%
  \rotatebox{45}{$\blacksquare$}}}}

%%%%%%%%%%%%%%%%%%%%%%%%%%%%%%%%%%%%%%%%%%%%%%%%%%%%%%%%%%%%%%%%%%%%%%%%
% Incipit
%%%%%%%%%%%%%%%%%%%%%%%%%%%%%%%%%%%%%%%%%%%%%%%%%%%%%%%%%%%%%%%%%%%%%%%%

\title{On Language Modeling and its Applications for Contemporary and Historical French}
%\subtitle{A minimal, modern \LaTeX{} package for typesetting your thesis}
\author{Pedro Ortiz Suarez}

\begin{document}

\frontmatter
\selectlanguage{french}

\begin{titlepage}

	\begin{center}
		\begin{minipage}[t]{0.4\textwidth}
			\includegraphics[height=2cm]{static/media/sorbonne}
		\end{minipage}%
		\hfill
		\begin{minipage}[t]{0.4\textwidth}
			\hfill
			\includegraphics[height=2cm]{static/media/inria}
		\end{minipage}

		\vspace{0.2cm}
		\LARGE \textsc{Sorbonne Université Université}\\
		\vspace{0.2cm}
		\normalsize \textsc{Ecole Doctorale Informatique, Télécommunications et Electronique} - ED130\\
		\vspace{0.2cm}
		\textsc{Inria de Paris / Équipe ALMAnaCH}\\

		\vspace{0.4cm}
		\Large \textsc{Thèse de doctorat}\\
		\normalsize Discipline : Informatique\\
		\vspace{0.4cm}
		\normalsize Présentée par \\
		\LARGE \textbf{Jean-Pierre \textsc{Dupont}}\\
		\vspace{0.4cm}
		\normalsize Dirigée par \\
		\Large \textbf{Marie \textsc{Curie}} et \textbf{Louis \textsc{Pasteur}}\\
		\vspace{0.4cm}
		\normalsize Pour obtenir le grade universitaire de\\
		\Large \textsc{Docteur} de \textsc{Sorbonne Université}

		\hrulefill\\[0.2cm]

		{\Large  \textbf{On Modern NLP for Extraterrestrial Languages and Scripts}}\\[0.1cm]

		\hrulefill\\

		\vspace{0.cm}
		\normalsize Présentée et soutenue publiquement le 30 avril 2022 devant le jury composé de :\\
		\vspace{0.4cm}
		\begin{tabular*}{\linewidth}{l c r}
			Alonzo \textsc{Church} & Princeton University & Examinateur \\
			Christopher \textsc{Columbus} & Kingdom of Castile & Invited member\\
			Margaret \textsc{Hamilton} & University of Michigan & Rapporteur \\
			Emmy \textsc{Noether} & Georg-August-Universität Göttingen & Examinateur \\
			Laurent \textsc{Romary} & Inria - ALMAnaCH & Directeur\\
			Benoît \textsc{Sagot} & Inria - ALMAnaCH & Co-directeur\\
			Claude \textsc{Shannon} & MIT & Examinateur \\
			Alan \textsc{Turing} & Princeton University & Rapporteur \\
		\end{tabular*}
	\end{center}

\end{titlepage}

\selectlanguage{english}

\newpage
\null
\thispagestyle{empty}
\newpage

\begin{center}
  \textsc{Abstract}
\end{center}
%
\noindent
%

The field of Natural Language Processing has recently known a major paradigm shift that has led to significant improvements over the perceived capabilities of resulting systems. This shift, namely the advent of generative systems in the stead of predictive ones, has induced a profound change in the implicit objectives of language systems based on deep learning : where we used to aim at extracting relevant features from text utterances using self-supervision, we now try to maximize the generative performance of language models on tremendous volumes of diversified text samples.

In this thesis, we explore high-level properties of the features (or \textit{representations}) that are extracted by these language models, and we leverage these properties to improve language systems and to quantify their limitations. This work is two-fold, as we first focus on the learnings that result from representation analysis in trained language models, and we then proceed to suggest and implement novel inductive biases and training approaches based on these learnings.

We find that the intermediate representations of self-supervised language models are affected by several forms of biases. First, they suffer from data-inherent biases that can be traced back to socio-cultural considerations, as we demonstrate by probing geographical knowledge in these features. Moreover, we show that these representational spaces can be distorted by the particular nature of language, especially when the dimensionality of the feature space is small. We proceed to show that inductive biases such as self-attention can also induce similar distortions that happen regardless of the target modality. Hence, representation analysis helps us identify limitations that come from distinct aspects of language models, from training data to architecture.

Not only can the prism of representation learning help us identify limitations in language models, but it can also lead to substantial improvements for language systems, especially in terms of efficiency. Aware of the mechanisms that we identified, we propose alternatives to the classical next-token likelihood maximization approach. We design a novel differentiable layer that performs text segmentation to optimize tokenization along with the rest of the system, leading to efficient character-level modeling and robust models. We also implement a contrastive objective that simultaneously alleviates the representational bias induced by token frequency and the degeneration phenomenon by implicitly regularizing the latent spaces using in-batch samples. This objective yields substantial efficiency improvements and better performance. Finally, we explore a differentiable compression scheme for generative language models, paving the way towards memory-efficient attention mechanisms.

Overall, our work proves the relevance of representation analysis in the context of improving language models, both as a way to identify limitations to the classical approaches and architectures, and as a way to gather insights on how to overcome these limitations.


\tableofcontents

\mainmatter

\part[A good part]{%
  A good part\\
  %
  \vspace{1cm}
  %
  \begin{minipage}[l]{\textwidth}
    %
    \textnormal{%
      \normalsize
      %
      \begin{singlespace*}
        \onehalfspacing
        %
        You can also use parts in order to partition your great work
        into larger `chunks'. This involves some manual adjustments in
        terms of the layout, though.
      \end{singlespace*}
    }
  \end{minipage}
}

%%%%%%%%%%%%%%%%%%%%%%%%%%%%%%%%%%%%%%%%%%%%%%%%%%%%%%%%%%%%%%%%%%%%%%%%
\chapter{Introduction}
%%%%%%%%%%%%%%%%%%%%%%%%%%%%%%%%%%%%%%%%%%%%%%%%%%%%%%%%%%%%%%%%%%%%%%%%

\begin{center}
  \begin{minipage}{0.5\textwidth}
    \begin{small}
      In which the reasons for creating this package are laid bare for the
      whole world to see and we encounter some usage guidelines.
    \end{small}
  \end{minipage}
  \vspace{0.5cm}
\end{center}

\noindent This package contains a minimal, modern template for writing your
thesis. While originally meant to be used for a Ph.\,D.\ thesis, you can
equally well use it for your honour thesis, bachelor thesis, and so
on---some adjustments may be necessary, though.

%%%%%%%%%%%%%%%%%%%%%%%%%%%%%%%%%%%%%%%%%%%%%%%%%%%%%%%%%%%%%%%%%%%%%%%%
\section{Background}
%%%%%%%%%%%%%%%%%%%%%%%%%%%%%%%%%%%%%%%%%%%%%%%%%%%%%%%%%%%%%%%%%%%%%%%%

I was not satisfied with the available templates for \LaTeX{} and wanted
to heed the style advice given by people such as Robert
Bringhurst~\cite{Bringhurst12} or Edward R.\
Tufte~\cite{Tufte90,Tufte01}. While there \emph{are} some packages out
there that attempt to emulate these styles, I found them to be either
too bloated, too playful, or too constraining. This template attempts to
produce a beautiful look without having to resort to ay sort of hacks.
I hope you like it.

%%%%%%%%%%%%%%%%%%%%%%%%%%%%%%%%%%%%%%%%%%%%%%%%%%%%%%%%%%%%%%%%%%%%%%%%
\section{Motivation}
%%%%%%%%%%%%%%%%%%%%%%%%%%%%%%%%%%%%%%%%%%%%%%%%%%%%%%%%%%%%%%%%%%%%%%%%

The package tries to be easy to use. If you are satisfied with the
default settings, just add
%
\begin{verbatim}
\documentclass{mimosis}
\end{verbatim}
%
at the beginning of your document. This is sufficient to use the class.
It is possible to build your document using either \LaTeX|, \XeLaTeX, or
\LuaLaTeX. I personally prefer one of the latter two because they make
it easier to select proper fonts.

%%%%%%%%%%%%%%%%%%%%%%%%%%%%%%%%%%%%%%%%%%%%%%%%%%%%%%%%%%%%%%%%%%%%%%%%

Prior to using this template, the first thing you want to do is probably
a little bit of customisation. You can achieve quick changes in look and
feel by picking your own fonts. With the \verb|fontspec| package loaded
and  \XeLaTeX or \LuaLaTeX as your compiler, this is pretty simple:
%
\begin{verbatim}
\setmainfont{Your main font}
\setsansfont{Your sans-serif font}
\setmonofont{Your monospaced font}
\end{verbatim}
%
Make sure to select nice combinations of that are pleasing to
\emph{your} eyes---this is your document and it should reflect your own
style. Make sure to specify font names as they are provided by your
system. For instance, you might want to use the following combination:
%
\begin{verbatim}
\setmainfont{Baskerville}
\setsansfont[Scale=MatchLowercase]{IBM Plex Sans}
\setmonofont[Scale=MatchLowercase]{IBM Plex Mono}
\end{verbatim}
%
%
You can also remove the \verb|Scale| directive, but I find that most
fonts pair better if they are adjusted in size a little bit. Experiment
with it until you finds a combination that you enjoy.

%%%%%%%%%%%%%%%%%%%%%%%%%%%%%%%%%%%%%%%%%%%%%%%%%%%%%%%%%%%%%%%%%%%%%%%%

%%%%%%%%%%%%%%%%%%%%%%%%%%%%%%%%%%%%%%%%%%%%%%%%%%%%%%%%%%%%%%%%%%%%%%%%
\begin{table}
  \centering
  \begin{tabular}{ll}
    \toprule
    \textbf{Package}      & \textbf{Purpose}\\
    \midrule
      \texttt{amsmath}          & Basic mathematical typography\\
      \texttt{amsthm}           & Basic mathematical environments for proofs etc.\\
      \texttt{booktabs}         & Typographically light rules for tables\\
      \texttt{bookmarks}        & Bookmarks in the resulting PDF\\
      \texttt{dsfont}           & Double-stroke font for mathematical concepts\\
      \texttt{graphicx}         & Graphics\\
      \texttt{hyperref}         & Hyperlinks\\
      \texttt{multirow}         & Permits table content to span multiple rows or columns\\ 
      \texttt{paralist}         & Paragraph~(`in-line') lists and compact enumerations\\
      \texttt{scrlayer-scrpage} & Page headings\\
      \texttt{setspace}         & Line spacing\\
      \texttt{siunitx}          & Proper typesetting of units\\
      \texttt{subcaption} & Proper sub-captions for figures\\
    \bottomrule
  \end{tabular}
  \caption{%
    A list of the most relevant packages required~(and automatically imported) by this template.
  }
  \label{tab:Packages}
\end{table}
%%%%%%%%%%%%%%%%%%%%%%%%%%%%%%%%%%%%%%%%%%%%%%%%%%%%%%%%%%%%%%%%%%%%%%%%

The template automatically imports numerous convenience packages that
aid in your typesetting process. \autoref{tab:Packages} lists the
most important ones. Let's briefly discuss some examples below. Please
refer to the source code for more demonstrations.

%%%%%%%%%%%%%%%%%%%%%%%%%%%%%%%%%%%%%%%%%%%%%%%%%%%%%%%%%%%%%%%%%%%%%%%%

Along with the standard environments, this template offers
\verb|paralist| for lists within paragraphs.
%
Here's a quick example: The American constitution speaks, among others, of
%
\begin{inparaenum}[(i)]
  \item life
  \item liberty
  \item the pursuit of happiness.
\end{inparaenum}
%
These should be added in equal measure to your own conduct. To typeset
units correctly, use the \verb|siunitx| package. For example, you might
want to restrict your daily intake of liberty to \SI{750}{\milli\gram}.

Likewise, as a small pet peeve of mine, I offer specific operators for
\emph{ordinals}. Use \verb|\th| to typeset things like July~4\th
correctly. Or, if you are referring to the 2\nd edition of a book,
please use \verb|\nd|. Likewise, if you came in 3\rd in a marathon, use
\verb|\rd|. This is my 1\st rule.

\section{Scope}
What we wanted to do.

%%%%%%%%%%%%%%%%%%%%%%%%%%%%%%%%%%%%%%%%%%%%%%%%%%%%%%%%%%%%%%%%%%%%%%%%
\section{Contributions}
%%%%%%%%%%%%%%%%%%%%%%%%%%%%%%%%%%%%%%%%%%%%%%%%%%%%%%%%%%%%%%%%%%%%%%%%

Since this class heavily relies on the \verb|scrbook| class, you can use
\emph{their} styling commands in order to change the look of things. For
example, if you want to change the text in sections to \textbf{bold} you
can just use
%
\begin{verbatim}
  \setkomafont{sectioning}{\normalfont\bfseries}
\end{verbatim}
%
at the end of the document preamble---you don't have to modify the class
file for this. Please consult the source code for more information.



% This ensures that the subsequent sections are being included as root
% items in the bookmark structure of your PDF reader.
\bookmarksetup{startatroot}
\backmatter

\begingroup
\let\clearpage\relax
\glsaddall
\printglossary[type=\acronymtype]
\newpage
\printglossary
\endgroup

\printindex

% Entries for the entire Anthology, followed by custom entries
\bibliography{anthology, thesis}
\bibliographystyle{acl_natbib}

\end{document}

